\subsection{Forest Disturbance}

In Kenya, approximately 75\% of all greenhouse gas (GHG) emissions related to forestry result not from large-scale timber harvesting, but from small-scale forest disturbances \parencite{pearsonGreenhouseGasEmissions2017}.
Consequently, the direct conversion of forested land to permanent settlements partly captured in H1 and H2 only represents less than a quarter of the total forest loss in Kenya.
To fully capture the extent of forest loss, it is necessary to consider a spillover effect, driven by the demand for forest products by settlers of nearby schemes and the access to those areas by new infrastructural developments.

The main mechanism for the disturbances is the extraction of Non-Timber Forest Products (NTFPs), specifically fuelwood and charcoal.
In a report to the United Nations Framework Convention on Climate Change (UNFCCC), the Government of Kenya estimated that approximately 70\% of the non-electric energy generation relies on fuelwood and charcoal.
This results in these resources being the primary driver of wood removal since 1990, the beginning of recording the wood removal sources \parencite{governmentofkenyaSecondNationalCommunication2015}.
Adding to this, \cite{hosonumaAssessmentDeforestationForest2012} estimated that fuelwood and charcoal production accounted for 58\% of forest degradation in a group of 33 countries, including Kenya by analyzing the data of 46 tropical and subtropical countries.

Following resettlement, smallholder farms often lack access to financial resources and market infrastructure, making them reliant on nearby forests.
\cite{naughton-trevesLessonsParksPoverty2011} found that poor households in Uganda relied on forests for subsistence, particularly during periods of crisis, when woodfuel is unaffordable or nearby communal resources are depleted.
For Kenya, \cite{geistProximateCausesUnderlying2002} found that farmers relying on permanent cultivation for subsistence are three times more likely to expand their agricultural plots into surrounding areas, compared to commercial farming.

This is confirmed for the Kenyan Mau Forest, where empirical evidence shows that one determinant of forest disturbance is the spatial proximity to farms.
Using the NDMI, as it is highly sensitive to smaller changes in tropical wet forests compared to the NDVI, \cite{brandtSustainableIntensificationDairy2018} observe disturbances on forest edges.
They confirm the harvesting of fuelwood and subsequent cattle grazing mainly by smallholder farms as the main driver of this degradation.
The two factors often occur in combination, as the removal of fuelwood leads to a reduction in tree canopy, which in turn created grazing areas, where the cattle hinders the regeneration of the forest \parencite{brandtSustainableIntensificationDairy2018}.

Additionally, the implementation of settlement schemes has led to the construction of infrastructure.
\cite{geistProximateCausesUnderlying2002} compile 93 studies showing the construction of roads is considered another large driver of tropical deforestation with a strong negative effect on local vegetation density \parencite{pfaffWhatDrivesDeforestation1999}.
Although the construction is driven by the need to connect the new settlements, it also creates access to previously remote areas, which \cite{oliveiraLandUseAllocationProtects2007} found to be a strong driver of deforestation in the Brazilian Amazon.

Looking at Kenya's history and the political economy, the argumentation of \cite{sanfordDemocratizationElectionsPublic2023} to incentivize this type of behavior can also be applied.
To reward political support, the government does not have to clear a forest and resettle people, but can instead provide access to forest resources for example by reducing the level of protection of a forest area or by suspending the enforcement.
This is also documented to some extent by the report by \cite{ndunguReportComissionInquiry2004}.

Combining the factor that the main driver of forest disturbance is the extraction of NTFPs with the fact that, the reliance on these resources by settlers,
the proximity of farms to forests, and the access to previously remote areas, this could lead to a spillover effect, which is captured in the third hypothesis:

Hypothesis 3 (H3): Settlement schemes generate a negative spillover effect, which leads to a higher rate of forest disturbance and deforestation in forested areas close to the scheme boundaries.